% !TEX root = ../main.tex
\chapter{Settings}\label{ch:settings}
Settings has sixteen pages, listed down its left side. Fourteen are \emph{project} settings: they are
stored in the project's database and apply to everyone who opens that database. Two, Display and
Keyboard \& behaviour, are \emph{personal}: they are stored in this browser and change only how the
site looks and responds to you. The table says what each page holds; the section after it explains
the two kinds, and the audit log that records every change.

On a project page your edits are a draft until you press \ui{Save}; \ui{Discard} drops them, and
leaving the page with unsaved edits asks first. A personal page applies each change at once.

\begin{longtable}{@{}>{\raggedright\arraybackslash}p{0.2\linewidth}>{\raggedright\arraybackslash}p{0.22\linewidth}>{\raggedright\arraybackslash}p{0.52\linewidth}@{}}
\toprule
\bfseries Page & \bfseries Applies to & \bfseries What it controls\\
\midrule\endhead
\pageanchor{settings.datasets}Datasets & the project; stored in its database & The registered
\term{recordings}{recording}, grouped by \term{dataset}{dataset}. Per dataset: its identity (display
name, species, organism id, experiment date, condition, notes), its recording conditions (substrate,
electrode configuration, start time, time zone, noise floor, temperature, humidity), the unit its
samples are stored in, and facts read from the data (sampling rate, channels, duration). \ui{Import a
recording\dots} registers a new file. The \term{held-out}{held-out} lock is set here; turning it off
needs the recording's name typed exactly. The noise floor is read by the Library's view filter
(\code{Working/library/view_filter.py}).\\
\pageanchor{settings.channels-events}Channels \& events & the project; stored in its database &
Per-channel metadata of one recording (status, ground, noise floor and the like) and its timed
\emph{event log}: marked moments or spans, each with a kind and an effect, plus your own event
kinds. Training-window pools read the added and removed events and their effects
(\code{Working/training/pool.py}).\\
\pageanchor{settings.vocabulary}Vocabulary & the project; stored in its database & The words a
person judges with: the five \term{verdicts}{verdict} (fixed) with their keys and how many
annotations and adjudications use each, the optional \term{classes}{class} for multi-class models,
and the morphology \term{tags}{tag} for Library entries. Tags can be added, renamed and merged.\\
\pageanchor{settings.nulls}Nulls & the project; stored in its database & Which
\term{null}{null} each kind of analysis uses (the method, from those the surrogate block offers, and
how many draws), the significance level, the correction for many channels, the block length of a
block shuffle, and whether null draws count toward the local time limits. Seed searches and an
Analyse run with its null switched on read the method and draws
(\code{Working/discovery/seeded_search.py}, \code{null_from_settings}).\\
\pageanchor{settings.analysis-defaults}Analysis defaults & the project; stored in its database &
The rule for when two events are the same (overlap and onset tolerance, read by
\code{Working/discovery/matching.py}), window-label containment, the seed-search scale bank and
exclusion zone, the cross-channel classification thresholds (\code{Working/cross_channel.py}), the
frequency bands a fan-out over bands uses (\code{Working/run_groups.py}), the reference for the
Library's rose plot, and the artifact flags shown in Review. It also shows the step cache's size and
can clear it.\\
\pageanchor{settings.compute-hpc}Compute \& HPC & the project; stored in its database & This
machine (detected cores and memory, how many local jobs run at once), the \emph{local limit} per
workspace above which work is written as a cluster job instead of run here, the cluster profiles
and the script preview, and how often to remind you about a running cluster job. The Analyse and
Discovery limits are read by the bridge and the core (\code{webui/server/chain.py},
\code{Working/discovery/fanout.py}); \unverified{no code reading the cluster profiles was found:
the job writers use the constants in \code{Working/config.py}}.\\
\pageanchor{settings.blocks}Blocks & the project; stored in its database & Every registered
\term{block}{block}, read from \code{Adapters/} when the bridge starts: its types, whether it is
known to be broken, and which templates use it. Each has an on/off switch;
\unverified{no code outside the Settings pages was found that reads the switch}.\\
\pageanchor{settings.review-queues}Review queues & the project; stored in its database & Defaults
for new \term{review queues}{review-queue} by source (order, cap, blind or not), batch verdicts on
clusters (cohesion limit, largest batch, undo), and promotion of a seed to the Library.
\unverified{Which of these the queue builder reads was not traced.}\\
\pageanchor{settings.models-registration}Models \& registration & the project; stored in its
database & Defaults for new training jobs: the evaluation split (test portion,
\term{blocked split}{blocked-split}, gap), label arms, repeats with different seeds, early
stopping, calibration, and the checks a \term{model}{model} must pass to be registered. Lists the
registered models. \unverified{Which defaults the training code reads was not traced.}\\
\pageanchor{settings.library-groupings}Library groupings & the project; stored in its database &
Defaults for new \term{groupings}{grouping}: the unit and basis, when a motif is omitted, when
events form one sequence, the feature bases, and whether hand edits are re-applied after a regroup.
\unverified{Which defaults the grouping engine reads was not traced.}\\
\pageanchor{settings.storage-backups}Storage \& backups & the project; stored in its database &
The folders the app reads and writes, each with a \ui{scan} that finds files to register
(\cref{sec:extending-data}); the naming convention for files stages write; database backups
(where, how many to keep, \ui{Back up now}, restore); and the free-space warning.\\
\pageanchor{settings.export}Export & the project; stored in its database & What an export carries:
the fields every export includes, the formats and contents for motifs and families, window sets,
templates, models and run reports, and the reproducibility bundle.\\
\pageanchor{settings.audit-log}Audit log & the project; read only & Every recorded change, newest
first, filterable by kind, with a link to where it happened and a CSV download. See below.\\
\pageanchor{settings.about}About & the project; read only & What this installation runs: the
code version, database schema, number of blocks and known-broken ones, the Python packages, the
mode (sandbox or project) and the database path, with a button that copies these as diagnostics.\\
\pageanchor{settings.display}Display & this browser only & Theme and density, how times,
amplitudes and sample indices are written, colours for roles and families, how many items a view
shows per page, and the figure profile run reports and figure exports use.\\
\pageanchor{settings.keyboard}Keyboard \& behaviour & this browser only & The keyboard shortcuts,
each rebindable (a clash with another key is shown), and how adjudication behaves in Review. The
verdict and class keys are part of the vocabulary and are set on Vocabulary.\\
\bottomrule
\end{longtable}

\screenshot{pages/settings/settings.datasets--default}{Settings opens on Datasets. The left rail
groups the fourteen project pages above the two personal ones.}

\section{Project and personal settings}\label{sec:settings-scopes}
\begin{plainterms}[Two kinds of setting]
Think of a shared lab notebook and your own desk. What is written in the notebook (which recordings
exist, what a verdict means, how long a run may take on the laptop) is the same for everyone who
opens it. How you arrange your desk (light or dark, which key does what) is yours alone.
\end{plainterms}

A \textbf{project} setting is a row of the \code{settings} table (\code{page}, \code{key}, a value
as JSON, when, by whom), read with \code{GET /api/settings/{page}} and saved with
\code{PUT /api/settings/{page}} (\code{webui/server/registration.py}). It lives in the project's
database, so it travels with that database and applies to every page and every person using it.
The six identity fields of a dataset are the exception: they are stored in the \code{datasets}
table, keyed by source file. A setting reaches a run only through the code that reads it: the
local limit decides where a run goes, the null settings decide which null a run carries, the
matching rule decides which detections count as the same event. The app does not store a copy of
every setting with every run; \unverified{whether a given result records the setting values it used
depends on the code that produced it}, and \cref{sec:architecture-reproducibility} says what a run
does record.

In the bridge's sandbox mode (\cref{sec:architecture-sandbox}) a save writes to the sandbox's copy
of the database, so it is gone when the sandbox is.

A \textbf{personal} setting never reaches the server. Display and Keyboard \& behaviour are kept in
this browser's \code{localStorage} (\code{webui/client/src/settings/store.ts}), so they differ
between browsers and between people, and a private window starts from the defaults.

\section{The audit log}\label{sec:settings-audit}
Every project save appends an entry to the audit log naming the page and each changed value, old
and new. So do turning the held-out lock off, registering or unregistering a file, declaring a
recording's unit, backups, and Library changes such as hand edits. The \code{audit_log} table
(\code{Working/database/schema.py}) is append-only: nothing in the code updates or deletes a row.
The Audit log page reads it through \code{GET /api/audit}, filters it by kind, links each entry to
the page it came from, and downloads it as CSV (\code{GET /api/audit.csv}).

\screenshot{pages/settings/settings.audit-log--all}{The audit log, newest first. Each entry says
what changed, where, and links back to that page.}
